\section{Results}
\label{sec:results}

\subsection{Temporal node-risk prediction}

We first examine whether short temporal histories improve the prediction
of next-interval node risk. Table~\ref{tab:model-quality} compares the
temporal histogram-gradient-boosting model with its current-state
counterpart on the final held-out node transitions.

\begin{table}[t]
\centering
\caption{Node-risk prediction on the final test transitions.}
\label{tab:model-quality}
\begin{tabular}{lrr}
\toprule
Metric &
Current-state HGB &
Temporal HGB \\
\midrule
ROC-AUC $\uparrow$ &
0.996944 &
\textbf{0.998248} \\

PR-AUC $\uparrow$ &
0.984894 &
\textbf{0.990308} \\

Brier score $\downarrow$ &
0.009858 &
\textbf{0.008578} \\
\bottomrule
\end{tabular}
\end{table}

Both models achieve high discrimination, but the temporal representation
improves every common test metric. The increase in PR-AUC is particularly
relevant because unsafe transitions form the minority class. The lower
Brier score further shows that temporal history improves the quality of
the predicted probabilities, although it does not identify a separate probability-calibration procedure.
For the temporal model, the final test log loss is 0.030480 and the
expected calibration error is 0.000970. Table~\ref{tab:operating-points} compares the learned temporal screens
with deterministic guardrail rules at closely matched acceptance levels.
The default temporal operating point accepts 82.35\% of node transitions,
which is close to the 82.11\% acceptance of the Q95 rule. Its unsafe rate
among accepted nodes is 0.0571\%, compared with 0.1343\% for Q95, while
retaining a slightly larger share of actually safe nodes. The same
pattern appears at the stricter operating point: temporal screening
achieves a lower unsafe rate than the deterministic 0.99 boundary at
nearly identical acceptance.

\begin{table*}[t]
\centering
\caption{Test-set operating-point comparison. Unsafe rate is measured
among accepted node transitions; safe retention is the fraction of
actually safe transitions retained by the rule.}
\label{tab:operating-points}
\begin{tabular}{llrrr}
\toprule
Screen &
Threshold or cutoff &
Acceptance (\%) &
Unsafe among accepted (\%) &
Safe retention (\%) \\
\midrule
Temporal availability &
$0.011362695$ &
82.349 &
\textbf{0.0571} &
\textbf{93.287} \\

Q95 guardrail &
$1.000$ &
82.106 &
0.1343 &
92.939 \\

Temporal safety &
$0.002059431$ &
75.714 &
\textbf{0.0236} &
85.799 \\

Boundary guardrail &
$0.990$ &
75.819 &
0.0617 &
85.885 \\

Strict Q95 guardrail &
$0.975$ &
62.960 &
0.0311 &
71.340 \\
\bottomrule
\end{tabular}
\end{table*}

These results establish that temporal history contributes information
that is not captured by current utilization or a fixed projected-headroom
rule. The default temporal threshold is therefore used in the complete
scheduler, while the deterministic Q95 rule remains as the final
pre-commit verification stage.

\subsection{Replay comparison on the natural workload}

Table~\ref{tab:natural-results} presents the main comparison on the
natural test workload of 11,629 events. The observed-source control
illustrates the difficulty of retaining the native placements: 25.875\%
of events remain on policy-inadmissible sources, 19.512\% exceed the
temporal-risk threshold, and 19.451\% fail the Q95 condition. Enforcing policy alone removes policy violations but leaves substantial
operational exposure. Policy-only-fair records temporal and Q95
violation rates of 13.011\% and 13.019\%, respectively. Conversely,
removing policy from the temporal scheduler preserves low operational
risk but leaves 25.411\% policy violations. These two controls show that
compliance and operational suitability address different requirements
and neither can replace the other.

\begin{table*}[t]
\centering
\caption{Replay results on the natural test workload
(11,629 events). Rates and Top-10 are percentages; Jain is unitless. Risk means temporal-screen violation. Unsafe is the native next-state proxy. Short names refer to the controls
in Section~\ref{sec:experimental}.}
\label{tab:natural-results}

\begin{tabular}{lrrrrrrr}
\toprule
Method &
Migr. &
Policy &
Risk &
Q95 &
Unsafe &
Jain &
Top-10 \\
\midrule
Observed source &
0.000 &
25.875 &
19.512 &
19.451 &
12.340 &
-- &
-- \\

Random &
38.198 &
0.000 &
1.392 &
0.000 &
0.190 &
0.191 &
1.353 \\

Least-loaded &
38.198 &
0.000 &
0.920 &
0.000 &
0.155 &
0.009 &
13.282 \\

Policy-only-fair &
25.875 &
0.000 &
13.011 &
13.019 &
7.215 &
0.189 &
0.665 \\

Policy-Q95-fair &
38.198 &
0.000 &
0.980 &
0.000 &
0.155 &
0.291 &
0.450 \\

Boundary-0.99 &
43.323 &
0.000 &
0.404 &
0.000 &
0.095 &
\textbf{0.332} &
\textbf{0.397} \\

Current-HGB &
42.850 &
0.000 &
0.370 &
0.000 &
0.069 &
0.329 &
0.401 \\

No policy &
20.741 &
25.411 &
0.000 &
0.000 &
0.060 &
0.190 &
0.415 \\

No projection &
38.189 &
0.000 &
0.000 &
1.144 &
\textbf{0.043} &
0.290 &
0.450 \\

\cage &
39.118 &
\textbf{0.000} &
\textbf{0.000} &
\textbf{0.000} &
\textbf{0.043} &
0.299 &
0.440 \\
\bottomrule
\end{tabular}

\end{table*}

The comparison with the current-state learned scheduler changes the
feature set and the associated admission threshold. Current-HGB migrates 42.850\%
of events and leaves a temporal-violation rate of 0.370\%. \cage reduces
migration to 39.118\%, removes the remaining temporal violations, and
reduces the next-state unsafe rate from 0.069\% to 0.043\%. The comparison with Policy-Q95-fair isolates the contribution of learned
temporal screening. \cage increases migration by only 0.920 percentage
points, while reducing the temporal-violation rate from 0.980\% to zero
and the next-state unsafe rate from 0.155\% to 0.043\%. Finally, the no-projection ablation isolates the pre-commit guardrail.
No projection already eliminates policy and temporal
violations, but 1.144\% of its final selections fail the Q95 condition.
Adding the pre-commit projection increases migration by 0.929 percentage
points and removes these violations without increasing the next-state
unsafe rate. The least-loaded baseline also demonstrates the importance of
fairness-aware selection. Although it satisfies policy and Q95
conditions, its migrated-destination Jain index is only 0.009 and its
ten most-used nodes receive 13.282\% of all migrations. Under the same
policy and Q95 conditions, Policy-Q95-fair increases Jain fairness to
0.291 and reduces the top-ten share to 0.450\%.

\subsection{Service-balanced workload}

The service-balanced workload tests whether the main observations depend
on the event concentration of a small number of services.
Table~\ref{tab:balanced-results} reports the results for 3,787 test
events.

\begin{table*}[t]
\centering
\caption{Replay results on the service-balanced test workload
(3,787 events). Rates and Top-10 are percentages; Jain is unitless. Risk means temporal-screen violation.}
\label{tab:balanced-results}

\begin{tabular}{lrrrrrrr}
\toprule
Method &
Migr. &
Policy &
Risk &
Q95 &
Unsafe &
Jain &
Top-10 \\
\midrule
Observed source &
0.000 &
24.611 &
21.917 &
21.310 &
13.890 &
-- &
-- \\

Random &
38.791 &
0.000 &
1.863 &
0.000 &
0.334 &
0.094 &
2.155 \\

Least-loaded &
38.791 &
0.000 &
1.400 &
0.000 &
0.290 &
0.002 &
39.959 \\

Policy-only-fair &
24.611 &
0.000 &
14.787 &
14.180 &
8.635 &
0.073 &
1.073 \\

Policy-Q95-fair &
38.791 &
0.000 &
1.400 &
0.000 &
0.290 &
0.115 &
0.681 \\

Boundary-0.99 &
43.623 &
0.000 &
0.528 &
0.000 &
0.158 &
\textbf{0.130} &
\textbf{0.605} \\

Current-HGB &
43.200 &
0.000 &
0.502 &
0.000 &
0.158 &
0.129 &
0.611 \\

No policy &
23.105 &
24.399 &
0.000 &
0.000 &
0.132 &
0.069 &
1.143 \\

No projection &
39.398 &
0.000 &
0.000 &
0.792 &
\textbf{0.106} &
0.117 &
0.670 \\

\cage &
40.190 &
\textbf{0.000} &
\textbf{0.000} &
\textbf{0.000} &
\textbf{0.106} &
0.120 &
0.657 \\
\bottomrule
\end{tabular}

\end{table*}

The balanced workload reproduces the natural-workload pattern. \cage
passes all configured policy, temporal-risk, and Q95 checks with a migration
rate of 40.190\% and a next-state unsafe rate of 0.106\%.
Current-HGB migrates more events, at 43.200\%, while leaving
0.502\% temporal violations. Policy-Q95-fair migrates slightly fewer
events, at 38.791\%, but its temporal-violation and next-state unsafe
rates increase to 1.400\% and 0.290\%, respectively. The fairness comparison becomes even more pronounced in this workload.
Least-loaded concentrates 39.959\% of migrations on its ten
most-used destinations and achieves a Jain index of only 0.002.
Policy-Q95-fair reduces the top-ten share to 0.681\% and raises the Jain
index to 0.115. \cage retains a similarly distributed assignment pattern
while adding temporal screening and pre-commit verification.



\subsection{Shortlist-size sensitivity}

Table~\ref{tab:shortlist-sensitivity} reports the candidate-selection
microbenchmark for fairness-aware shortlists. Full fairness-aware
candidate evaluation requires 72.930~ms on average. Restricting the
initial candidate set sharply reduces this cost.

\begin{table}[t]
\centering
\caption{Mean fairness-aware candidate-selection latency on the natural
workload. Feature construction and model inference are excluded from
this isolated shortlist microbenchmark.}
\label{tab:shortlist-sensitivity}
\begin{tabular}{lr}
\toprule
Candidate setting & Mean latency (ms) \\
\midrule
$k=128$ & 4.182 \\
$k=256$ & 4.724 \\
$k=512$ & 5.754 \\
$k=1024$ & 8.100 \\
$k=2048$ & 12.517 \\
Full candidate set & 72.930 \\
\bottomrule
\end{tabular}
\end{table}

The $k=256$ configuration reproduces the full fairness-aware selections,
requires no search expansion, and preserves a minimum accepted-candidate
slack of 130 nodes on the natural workload and 229 nodes on the balanced
workload. Its isolated shortlist latency has a P95 of 5.937~ms and a P99
of 7.576~ms. The smaller $k=128$ setting is faster, but its minimum
natural-workload slack falls to two candidates. \cage therefore uses
$k=256$ as the default configuration.

\subsection{Paired timestamp-cluster comparisons}

Table~\ref{tab:bootstrap-results} reports the principal paired
timestamp-cluster bootstrap comparisons. Differences are calculated as
\cage minus the corresponding baseline. For migration and violation
rates, an interval below zero favours \cage.

\begin{table*}[t]
\centering
\caption{Selected paired timestamp-cluster bootstrap comparisons.
Intervals are 95\% percentile confidence intervals based on 20,000
bootstrap samples. Values are proportions.}
\label{tab:bootstrap-results}
\begin{tabular}{llrrr}
\toprule
Workload &
Baseline and metric &
Difference &
CI low &
CI high \\
\midrule
Natural &
Current-HGB: migration &
-0.037320 &
-0.040806 &
-0.033591 \\

Natural &
Current-HGB: temporal violation &
-0.003698 &
-0.005049 &
-0.002499 \\

Natural &
Current-HGB: next-state unsafe &
-0.000258 &
-0.000574 &
0.000000 \\

Natural &
Policy-Q95: migration &
0.009201 &
0.007527 &
0.010939 \\

Natural &
Policy-Q95: temporal violation &
-0.009803 &
-0.011684 &
-0.008022 \\

Natural &
Policy-Q95: next-state unsafe &
-0.001118 &
-0.001751 &
-0.000571 \\

Natural &
No projection: migration &
0.009287 &
0.007350 &
0.011482 \\

Natural &
No projection: Q95 violation &
-0.011437 &
-0.013933 &
-0.009146 \\
\midrule
Balanced &
Current-HGB: migration &
-0.030103 &
-0.035847 &
-0.024455 \\

Balanced &
Current-HGB: temporal violation &
-0.005017 &
-0.007339 &
-0.002936 \\

Balanced &
Policy-Q95: migration &
0.013995 &
0.010437 &
0.017647 \\

Balanced &
Policy-Q95: temporal violation &
-0.013995 &
-0.017647 &
-0.010437 \\

Balanced &
Policy-Q95: next-state unsafe &
-0.001848 &
-0.003484 &
-0.000528 \\

Balanced &
No projection: migration &
0.007922 &
0.005424 &
0.010571 \\

Balanced &
No projection: Q95 violation &
-0.007922 &
-0.010571 &
-0.005424 \\
\bottomrule
\end{tabular}
\end{table*}

Relative to the current-state learned scheduler, \cage achieves
significantly lower migration and temporal-violation rates on both
workloads. Relative to Policy-Q95-fair, it incurs a small and precisely
estimated migration increase while significantly reducing temporal
violations and the next-state unsafe rate. Relative to the no-projection
ablation, the additional migration is accompanied by the complete removal
of Q95 violations. The natural-workload next-state interval against Current-HGB includes
zero; this comparison does not establish a strict improvement in that
proxy at the reported confidence level. Temporal violations measure
agreement with the temporal model, while the migration difference is an
independent decision-count comparison.

\subsection{Robustness across policy overlays}

The integrated policy--temporal--projection scheduler is evaluated across
30 deterministic policy overlays. Together, these runs contain 462,480
decisions. Every selected destination satisfies the policy,
temporal-risk, and Q95 conditions. No destination conflict, migration
into another active source, fallback, or missing decision is observed.

\begin{table*}[t]
\centering
\caption{Robustness across 30 deterministic policy overlays. Rates are
percentages except for Jain fairness.}
\label{tab:seed-robustness}
\begin{tabular}{lrrrrrr}
\toprule
Workload &
Migr. mean &
Migr. P05 &
Migr. P95 &
Unsafe mean &
Jain mean &
Jain P05--P95 \\
\midrule
Natural &
35.233 &
30.150 &
44.803 &
0.0496 &
0.3155 &
0.2722--0.3948 \\

Service-balanced &
37.837 &
34.927 &
41.211 &
0.1056 &
0.1126 &
0.1039--0.1226 \\
\bottomrule
\end{tabular}
\end{table*}

For the natural workload, the mean migration rate is 35.233\%, with a
5th--95th percentile range of 30.150--44.803\%. The mean next-state
unsafe rate remains 0.0496\%. For the service-balanced workload, the
corresponding values are 37.837\% and 0.1056\%. Variation in migration
reflects changes in the policy-admissible sets produced by each overlay;
the event sequence, temporal model, risk threshold, and guardrail
parameters remain fixed.




\subsection{Online decision latency}
\label{sec:runtime}

Table~\ref{tab:runtime-results} reports the latency of the complete
in-memory online core. The measurement includes temporal feature
construction, all required model inference, policy filtering,
fairness-aware shortlisting, ranking, reservation updates, and
pre-commit verification.



\begin{table}[t]
\centering
\caption{Complete online-core latency in milliseconds.}
\label{tab:runtime-results}
\begin{tabular}{lrr}
\toprule
Statistic &
Natural &
Service-balanced \\
\midrule
Mean & 20.247 & 13.152 \\
Median & 17.710 & 12.990 \\
P95 & 32.081 & 16.229 \\
P99 & 57.513 & 17.431 \\
Maximum & 60.293 & 20.247 \\
\bottomrule
\end{tabular}
\end{table}

The natural workload requires more computation because it contains more
events per timestamp and produces a larger policy-constrained inference
superset. Even its maximum observed online-core latency remains small
relative to the 30-second control interval. The online feature implementation is also checked against the stored
offline snapshots. The maximum feature difference is
$2.22\times10^{-16}$, the maximum probability difference is zero, and
no online placement decision differs from its stored replay reference.

\subsection{Destination-capacity sensitivity}

The main scheduler permits one new destination reservation per node and
timestamp. Repeating the complete replay with capacity $C=2$ does not
change any migration indicator. For the service-balanced workload, all 3,787 destinations remain
identical to the $C=1$ decisions. For the natural workload, 11,617 of
11,629 selected nodes remain identical, corresponding to a destination
agreement of 99.897\%. The 12 changed destinations do not change the
migration pattern, destination fairness, or any admission outcome. No node receives two assignments at the same timestamp under the
$C=2$ setting. The observed event density therefore does not make the
default capacity-one constraint an active placement bottleneck.
